\section{Kết luận}

\subsection{Tóm tắt Kết quả}
Đồ án đã nghiên cứu và xây dựng thành công hệ thống phân vùng thực thể chẩn đoán bệnh trên lá cây nông nghiệp đặc sản lúa và cà phê tại Việt Nam. Kết quả đạt được so với các mục tiêu đề ra ban đầu được tóm tắt như sau:
\begin{itemize}
  \item Xây dựng và chuẩn hóa thành công tập dữ liệu gồm 7.149 ảnh thực địa sạch cho cả cây lúa và cây cà phê sau khi loại bỏ các ảnh trùng lặp và ảnh mờ nhòe bằng thuật toán mã băm MD5 kết hợp kiểm tra phương sai Laplacian. Tập dữ liệu chứa 8 lớp nhãn chi tiết với các vết bệnh được khoanh vùng tỉ mỉ ở cấp độ điểm ảnh.
  \item Huấn luyện và đánh giá đối chiếu thực nghiệm 5 kiến trúc phân vùng thực thể hiện đại bao gồm hai hướng tiếp cận chính là mạng thần kinh tích chập và mạng thị giác máy tính dựa trên cơ chế tập trung.
  \item Lựa chọn mô hình YOLO26-seg làm mô hình tối ưu cho triển khai thực tế. Bằng việc áp dụng phương pháp tinh chỉnh siêu tham số thông qua thuật toán dừng sớm bất đồng bộ kết hợp Ray Tune, hiệu năng của mô hình được cải thiện thêm hơn 2\% độ chính xác mAP@50 trên cả hai miền dữ liệu, đạt lần lượt 84,81\% ở lúa và 92,34\% ở cà phê.
  \item Tối ưu hóa suy luận bằng cách lượng tử hóa mô hình sang định dạng ONNX. Qua đó, mô hình PyTorch \texttt{.pt} gốc giữ được dung lượng siêu nhẹ (6,23 MB), còn phiên bản ONNX INT8 phục vụ suy luận CPU đạt thời gian xử lý cực nhanh (trung bình 19,85 ms với ảnh cà phê và 65,41 ms với ảnh lúa) với dung lượng đóng gói 24,9 MB.
  \item Xây dựng hoàn chỉnh hệ thống ứng dụng web đa thành phần, hỗ trợ nông dân tải hoặc chụp ảnh chẩn đoán vết bệnh theo thời gian thực, tích hợp cơ sở dữ liệu lịch sử và hệ thống khuyến nghị tri thức chuyên gia nông nghiệp Việt Nam.
  \item Đóng gói và triển khai thành công ứng dụng trên môi trường đám mây sử dụng cụm Kubernetes k3s trên máy ảo GCP, cấu hình Traefik Ingress và bảo mật HTTPS, đồng thời tự động hóa toàn bộ quy trình bằng hệ thống tích hợp và triển khai liên tục.
\end{itemize}

\subsection{Hạn chế}
Mặc dù đã đạt được nhiều kết quả khả quan về cả độ chính xác của mô hình lẫn tính khả thi của hệ thống triển khai, đồ án vẫn tồn tại một số hạn chế cần được thẳng thắn nhìn nhận để cải thiện trong tương lai:
\begin{itemize}
  \item Sự mất cân bằng giữa các lớp dữ liệu vẫn còn khá lớn. Điển hình là bệnh phấn trắng trên cây cà phê chỉ chiếm một tỷ lệ rất nhỏ trong tập dữ liệu thu thập, làm giảm khả năng học và nhận diện các vết bệnh này của mô hình so với các bệnh phổ biến như rỉ sắt hay sâu đục lá.
  \item Mô hình chẩn đoán vẫn có xu hướng dự đoán sai hoặc giảm độ chính xác khi gặp các hình ảnh thực tế có điều kiện ánh sáng cực đoan ngoài đồng ruộng như ánh nắng gắt trực diện tạo bóng lóa, hoặc ảnh chụp ngược sáng mạnh làm mất chi tiết màu sắc đặc trưng của vết bệnh.
  \item Hạ tầng triển khai trên đám mây hiện đang vận hành dưới dạng cụm đơn node. Điều này dẫn đến nguy cơ cao về điểm lỗi đơn, nghĩa là nếu máy chủ vật lý gặp sự cố, toàn bộ dịch vụ sẽ ngừng hoạt động.
  \item Hệ thống lưu trữ tệp tin hình ảnh người dùng tải lên và cơ sở dữ liệu lịch sử chẩn đoán chưa được cấu hình cơ chế sao lưu dự phòng phân tán và nhân bản dữ liệu tự động giữa nhiều vùng lưu trữ khác nhau.
  \item Trình kiểm thử khói trong quy trình tự động hóa chưa bao phủ hết các ca kiểm thử suy luận ảnh thực tế từ đầu đến cuối, dẫn đến việc khó phát hiện sớm các lỗi tích hợp mô hình phát sinh sau mỗi lần cập nhật mã nguồn backend.
\end{itemize}

\subsection{Hướng phát triển}
Nhằm khắc phục các hạn chế nêu trên và nâng cao hơn nữa giá trị thực tiễn của đồ án, nhóm đề xuất các hướng nghiên cứu và phát triển tiếp theo như sau:
\begin{itemize}
  \item Tiếp tục mở rộng quy mô tập dữ liệu thực địa, tập trung thu thập hình ảnh của các lớp bệnh thiểu số để giải quyết triệt để vấn đề mất cân bằng lớp. Nghiên cứu áp dụng các phương pháp sinh ảnh hiện đại nhằm tạo thêm các mẫu bệnh đa dạng trong các điều kiện ánh sáng khác nhau.
  \item Khảo sát và thử nghiệm các kiến trúc phân vùng thực thể thế hệ mới có khả năng học đặc trưng kết cấu tốt hơn, hoặc tích hợp thêm mô hình ngôn ngữ lớn thị giác để hỗ trợ nông dân tương tác hỏi đáp thông tin bệnh học một cách trực quan, sinh động.
  \item Nâng cấp hạ tầng đám mây lên cụm Kubernetes nhiều node phân tán, sử dụng các giải pháp lưu trữ bền vững có hỗ trợ cơ chế nhân bản dữ liệu tự động nhằm loại bỏ hoàn toàn điểm lỗi đơn và nâng cao khả năng chịu lỗi của hệ thống.
  \item Tối ưu hóa phân phối tải cho backend bằng cách thiết lập cơ chế tự động mở rộng quy mô dựa trên lưu lượng truy cập thực tế, đồng thời giới hạn chính sách chia sẻ tài nguyên nguồn gốc chéo để đảm bảo an toàn thông tin cho API backend.
  \item Bổ sung hệ thống giám sát và cảnh báo tự động để theo dõi liên tục hiệu năng hoạt động của mô hình suy luận và tài nguyên hệ thống trên môi trường production, giúp phát hiện sớm các hiện tượng lệch dữ liệu theo thời gian.
\end{itemize}
